\section{Problem Formulation}
\label{sec:problem}

We study \emph{open-ended, evidence-grounded long-form generation} tasks where the final output is an expert-written artifact, but the generative process is unobserved. Given only an artifact $y^{*}$ (e.g., a related-work section, financial analysis report, or legal judgment), our goal is to reconstruct a plausible evidence-seeking trajectory that yields it.  

\paragraph{Latent Expert Process.}
We formulate $y^*$ as the outcome of an unobserved expert workflow:
\begin{equation}
(x^*, \mathcal{E}^*, \tau^*) \mapsto y^*,
\end{equation}
where $x^*$ is the underlying task specification, $\mathcal{E}^{*}$ represents the supporting evidence consulted by the expert, and $\tau^*$ is the latent process transforming the task and evidence into the final artifact. Because these variables are unobserved, we rely on the assumption that $y^{*}$ serves as a lossy but highly informative compression of this latent process.  


\paragraph{Retrospective Reconstruction.}
Recognizing that the exact historical process $\tau^{*}$ is unknowable, we instead aim to reconstruct a plausible sequence:
\begin{equation}
\hat{x} = g_x(y^*), \;
\hat{\mathcal{E}} = g_e(y^*), \;
\hat{\tau} = g_\tau(y^*, \hat{x}, \hat{\mathcal{E}}),
\end{equation}
where $\hat{x}$ is the reconstructed task, $\hat{\mathcal{E}}$ is the recovered evidence set, and $\hat{\tau}$ is a reverse-engineered trace describing how an agent could derive an output consistent with $y^*$.
Because the exact historical process $\tau^{*}$ is unobserved, we treat $\hat{\tau}$ as an \emph{executable surrogate trajectory}: a tool-using workflow that is faithful to $y^{*}$ and useful for training, rather than a reconstruction of the expert's unrecorded trial-and-error history.

\paragraph{Desiderata.}
A high-quality reconstructed trace $\hat{\tau}$ should satisfy three properties:

\begin{itemize}[leftmargin=*, topsep=3pt, partopsep=0pt, itemsep=0pt, parsep=3pt] 
    \item Artifact Fidelity: Preserving the essential content and structural organization of $y^{*}$.
    \item Evidence Faithfulness: Grounding intermediate reasoning and final claims strictly within the recovered evidence $\hat{\mathcal{E}}$.
    \item Procedural Plausibility: Exhibiting a logical domain workflow (e.g., search, inspect, compare) rather than hallucinating leaps in logic.
\end{itemize}

% \paragraph{Training objective.}
% Given a collection of expert artifacts
% \begin{equation}
% \mathcal{D}_{\mathrm{art}} = \{y_i^*\}_{i=1}^{N},
% \end{equation}
% our goal is to construct a retrospective process-supervision dataset
% \begin{equation}
% \hat{\mathcal{D}}_{\mathrm{RPS}}
% =
% \{(\hat{x}_{ij}, \hat{\mathcal{E}}_{ij}, \hat{\tau}_{ij}, y_i^*)\},
% \end{equation}
% where each retained tuple contains a plausible reconstructed process for artifact $y_i^*$. Multiple traces may be constructed for a single artifact, and low-quality traces may be filtered out. 

% This dataset is then used to train models that generate evidence-grounded long-form outputs, either directly or through tool-using interaction.